% Generated draft: requires mathematical and source review.
\documentclass[11pt]{article}
\usepackage[margin=25mm]{geometry}
\usepackage{fontspec}
\setmainfont{Latin Modern Roman}
\usepackage{amsmath,amssymb,mathtools,unicode-math}
\setmathfont{Latin Modern Math}
\usepackage{enumitem,needspace,xurl,hyperref}
\setlength{\parindent}{0pt}
\setlength{\parskip}{5pt}
\setlength{\emergencystretch}{3em}
\newcommand{\N}{\mathbb{N}}
\newcommand{\Z}{\mathbb{Z}}
\newcommand{\Q}{\mathbb{Q}}
\newcommand{\R}{\mathbb{R}}
\newcommand{\C}{\mathbb{C}}
\newcommand{\F}{\mathbb{F}}
\newcommand{\A}{\mathbb{A}}
\newcommand{\PP}{\mathbb{P}}
\newcommand{\E}{\mathbb{E}}
\newcommand{\eps}{\varepsilon}
\newcommand{\dd}{\,\mathrm{d}}
\begin{document}
\section*{UnsolvedMath ID: 30000779}
\textbf{Source problem code:} catalogue entry
\subsection{Definitions and mathematical statement}
Let \(k\) be a number field, let \(p\) be a rational prime, and let \(\mathcal O_k\) be the ring of integers. Write \(S_p\) for the set of nonzero prime ideals of \(\mathcal O_k\) above \(p\). Let \(S\) be a finite set of nonzero prime ideals with \(S\cap S_p=\varnothing\), and let \(T\) be a set of nonzero prime ideals of Dirichlet density one; explicitly,
\[
\lim_{s\downarrow1}\frac{\sum_{\mathfrak q\in T}(\mathrm N\mathfrak q)^{-s}}{\log(1/(s-1))}=1.
\]
The restored question in the supplied record is whether, for every such \(k,p,S,T\), there exists a finite subset \(T_1\subseteq T\) such that
\[
X=\operatorname{Spec}(\mathcal O_k)\setminus(S\cup T_1)
\quad\text{is an \(\acute{e}tale\) \(K(\pi,1)\) for \(p\)}.
\]
Here the removed prime ideals are viewed as closed points. In the usual pro-\(p\) interpretation, if \(G=\pi_1^{\acute{e}t}(X,\bar x)(p)\) is the maximal pro-\(p\) quotient of the \(\acute{e}tale\) fundamental group, the required property is that the natural comparison maps
\[
H^i_{\mathrm{cont}}(G,\F_p)\longrightarrow H^i_{\acute{e}t}(X,\F_p)
\]
are isomorphisms in every degree \(i\geq0\). The precise \(K(\pi,1)\)-for-\(p\) and real-place conventions of the cited theorem should be confirmed before treating this interpretation as definitive.

Crucially, the question imposes neither \(\zeta_p\notin k\), where \(\zeta_p\) is a primitive \(p\)-th root of unity, nor \(p\nmid|\operatorname{Cl}(k)|\), where \(\operatorname{Cl}(k)\) is the ideal class group. Both auxiliary restrictions are to be removed. The original extraction calls this removal of condition \((*)\) from the theorem; the complete theorem and the text of \((*)\) are not reproduced.
\subsection{Short English statement}
Can finitely many primes from any Dirichlet-density-one set be added to \(S\), making the resulting arithmetic curve an \(\acute{e}tale\) \(K(\pi,1)\) for \(p\), without either the cyclotomic or class-number restriction?
\subsection{Sources}
\begin{itemize}
\item Local UnsolvedMath registry, permanent ID \(30000779\), source code catalogue entry, ``\(K(\pi,1)\) Theorems without Cyclotomic Class-Number Conditions.'' No individual proposer is supplied.
\item \emph{Algebraische Zahlentheorie} (2008), Oberwolfach Reports, DOI 10.4171/owr/2007/30, pp. 1741--1792: \url{https://doi.org/10.4171/owr/2007/30}. This is the sole supplied external reference.
\item The literature triage checked 2026-08-21 labels the conjecture open and says that the cited report does not establish a later proof or counterexample. This does not independently establish its current status.
\end{itemize}
\end{document}
